\section{Resultados}
\label{sec:resultados}

Salvo indicação, as medidas usam $h = 0{,}05$ (20 células por diâmetro), $np=3$,
e o resíduo de não escorregamento é $\max_k |\mathbf{u}(\mathbf{X}_k)|$, a ser
comparado com $\max|\mathbf{u}| \approx 1{,}8$ no domínio.

\subsection{Estabilidade: o maior \texorpdfstring{$\dt$}{dt}, forçamento defasado}

\pgfplotstabletypeset[
  % TODAS string: os valores divergidos estouram o parser numerico do TeX.
  columns/0/.style={column name={$\dt$}, string type},
  columns/1/.style={column name={resíduo}, string type},
  columns/2/.style={column name={$\max|\mathbf{u}|$}, string type},
  columns/3/.style={column name={veredito}, string type},
  every head row/.style={before row=\toprule, after row=\midrule},
  every last row/.style={after row=\bottomrule},
]{dados/varre-dt-defasado.dat}

\medskip
Número de passos \emph{fixo} (400) para todos os valores: instabilidade de
realimentação cresce por passo, e a tempo fixo cada $\dt$ teria um número
diferente de oportunidades de amplificação.

\paragraph{A conclusão é o contrário do que motivou o experimento.} O
despenhadeiro está entre $5\times 10^{-2}$ e $10^{-1}$, \emph{acima} do CFL
convectivo ($\approx 0{,}026$): a fronteira imersa não é o gargalo de
estabilidade. O limite útil é de \textbf{acurácia} --- em $\dt = 5\times10^{-2}$,
último estável, o resíduo é $63\%$ de $\max|\mathbf{u}|$, e o obstáculo deixou de
ser obstáculo.

\paragraph{Advertência de método.} A primeira versão deste varrimento foi até
$2\times10^{-2}$ e \textbf{não encontrou limite algum}. Uma tabela inteiramente
verde \emph{parece} resposta. Varrimento que não encontra a fronteira não mediu o
que se pensa que mediu.

\subsection{O erro do forçamento defasado escala com o número de Courant}

\pgfplotstabletypeset[
  columns/0/.style={column name={$h$}, fixed, precision=4},
  columns/1/.style={column name={cél./$D$}, fixed, precision=0},
  columns/2/.style={column name={$\dt$}, sci, precision=2},
  columns/3/.style={column name={resíduo}, sci, precision=3},
  every head row/.style={before row=\toprule, after row=\midrule},
  every last row/.style={after row=\bottomrule},
]{dados/refino-defasado.dat}

\medskip
Tempo físico fixo ($t = 0{,}2$). O resíduo é $O(\dt)$ \emph{puro} nas duas
malhas --- razões $2{,}04$ e $2{,}02$ --- e \textbf{sem piso}. E a malha fina dá
o dobro da grossa no mesmo $\dt$.

Juntando as duas observações, o resíduo escala com $\dt/h$: o número de Courant.
Faz sentido --- o erro da defasagem é $\dt\,\partial_t u$, e junto ao corpo
$\partial_t u \sim u^2/h$. \textbf{Refinar a malha sem corrigir a defasagem
piora o resíduo}, e é por isso que refinamento local só faz sentido \emph{depois}
da mudança de formulação.

\subsection{Rota pós-preditor e forçamento iterado}

\pgfplotstabletypeset[
  columns/0/.style={column name={iterações}, fixed, precision=0},
  columns/1/.style={column name={resíduo}, sci, precision=3},
  columns/2/.style={column name={ganho vs.\ defasado}, fixed, precision=1},
  every head row/.style={before row=\toprule, after row=\midrule},
  every last row/.style={after row=\bottomrule},
]{dados/iteracoes.dat}

\medskip
Com $\dt = 10^{-3}$ e $t = 0{,}2$. Uma aplicação corta o resíduo pela
\emph{metade}, não o zera --- consequência direta de
$\mathcal{S}\mathcal{I}\neq\mathcal{I}$, e a razão pela qual a literatura
itera~\cite{wang2008}. A convergência satura: dobrar de 5 para 10 iterações
compra $1{,}38\times$; de 10 para 20, $1{,}40\times$.

\subsection{\texorpdfstring{$\dt$}{dt} com a rota pós-preditor}

\pgfplotstabletypeset[
  columns/0/.style={column name={$\dt$}, sci, precision=2},
  columns/1/.style={column name={Uhlmann, 5 it.}, sci, precision=3},
  columns/2/.style={column name={defasado}, sci, precision=3},
  every head row/.style={before row=\toprule, after row=\midrule},
  every last row/.style={after row=\bottomrule},
]{dados/varre-dt-uhlmann.dat}

\medskip
Abaixo de $2\times10^{-3}$ o resíduo \textbf{acha um patamar} ($\approx
1{,}7\times10^{-3}$): é o piso que a formulação defasada nunca deixou aparecer,
porque a rampa $O(\dt)$ o encobria.

\paragraph{O ganho se converte em custo de máquina.} A rota pós-preditor em
$\dt = 2\times10^{-2}$ é tão acurada quanto a defasada em $\dt = 10^{-3}$:
\textbf{vinte vezes o passo para a mesma acurácia}.

\subsection{Refinamento: \texorpdfstring{$\Cd$}{Cd} contra a referência}

\pgfplotstabletypeset[
  columns/0/.style={column name={cél./$D$}, fixed, precision=0},
  columns/1/.style={column name={$h$}, fixed, precision=4},
  columns/2/.style={column name={$\Cd$}, fixed, precision=6},
  columns/3/.style={column name={$\Cl$}, fixed, precision=6},
  columns/4/.style={column name={resíduo}, sci, precision=3},
  every head row/.style={before row=\toprule, after row=\midrule},
  every last row/.style={after row=\bottomrule},
]{dados/refino-cd.dat}

\medskip
Rota pós-preditor com 5 iterações, $\dt = 5\times10^{-3}$, medidos no
\emph{mesmo} instante $t = 14$, com $\Cd$ estável na quinta casa decimal entre
quadros consecutivos.

\begin{center}
\begin{tabular}{lSS}
  \toprule
  & {20 cél./$D$} & {40 cél./$D$} \\
  \midrule
  erro em $\Cd$ & {$+2{,}62\%$} & {$+0{,}41\%$} \\
  \bottomrule
\end{tabular}
\end{center}

A razão entre os erros é $6{,}4$ para um refinamento de $2$, o que dá inclinação
de convergência $\approx 2{,}7$.

\paragraph{Esse valor merece desconfiança.} A literatura mede inclinação
$\approx 1{,}0$ para forçamento direto simples e $\approx 1{,}5$ para
\emph{multi-direct forcing}. Obter $2{,}7$ é quase o dobro da melhor taxa
publicada. Três explicações competem: (i) as iterações compram ordem; (ii) o erro
da malha grossa tinha componentes que se cancelavam; (iii) algum detalhe da
medida de $\Cd$ escala com $h$ de modo não previsto. \textbf{Dois pontos não
distinguem as três} --- daí a terceira resolução.

\paragraph{Observação que separa duas coisas.} A malha fina tem resíduo de não
escorregamento \emph{pior} ($8{,}3\times10^{-3}$ contra $4{,}0\times10^{-3}$) e
$\Cd$ \emph{melhor}. Resíduo e erro no arrasto não são a mesma quantidade: o
primeiro mede escorregamento nos marcadores; o segundo depende de quão bem o
escoamento em torno do corpo está resolvido.
